% Generated draft: requires mathematical and source review.
\documentclass[11pt]{article}
\usepackage[margin=25mm]{geometry}
\usepackage{fontspec}
\setmainfont{Latin Modern Roman}
\usepackage{amsmath,amssymb,mathtools,unicode-math}
\setmathfont{Latin Modern Math}
\usepackage{enumitem,needspace,xurl,hyperref}
\setlength{\parindent}{0pt}
\setlength{\parskip}{5pt}
\setlength{\emergencystretch}{3em}
\newcommand{\N}{\mathbb{N}}
\newcommand{\Z}{\mathbb{Z}}
\newcommand{\Q}{\mathbb{Q}}
\newcommand{\R}{\mathbb{R}}
\newcommand{\C}{\mathbb{C}}
\newcommand{\F}{\mathbb{F}}
\newcommand{\A}{\mathbb{A}}
\newcommand{\PP}{\mathbb{P}}
\newcommand{\E}{\mathbb{E}}
\newcommand{\eps}{\varepsilon}
\newcommand{\dd}{\,\mathrm{d}}
\begin{document}
\section*{UnsolvedMath ID: 3100030}
\textbf{Source problem code:} catalogue entry
\subsection{Definitions and mathematical statement}
The Fano plane is the $3$-uniform hypergraph with vertex set $V=\mathbb F_2^3\setminus\{0\}$ and edges the distinct triples $\{u,v,u+v\}$ for $u\ne v$. It has seven vertices and seven edges. In the standard symmetric adjacency-tensor normalization, let
\[
a_{ijk}=\begin{cases}1/2,&\{i,j,k\}\text{ is an edge},\\0,&\text{otherwise}.\end{cases}
\]
A homogeneous tensor eigenvalue is $\lambda\in\mathbb C$ for which some nonzero vector $x\in\mathbb C^7$ satisfies
\[
\sum_{\{i,j,k\}\in E}x_jx_k=\lambda x_i^2\qquad(i\in V),
\]
where in the equation for $i$ the sum ranges over its incident edges and $j,k$ are their other vertices. This is equivalently $\mathcal A x^2=\lambda x^{[2]}$.

Determine the homogeneous adjacency spectrum of the Fano plane, including the intended multiplicities. Under the complex resultant convention the characteristic polynomial is the monic normalization of the homogeneous resultant of these seven equations, and the spectrum is its multiset of roots. The supplied item does not define its homogeneous-adjacency normalization, real-versus-complex eigenvalue convention, or multiplicity convention, so these choices are explicit standard context rather than a checked source definition. This is not the ordinary incidence-graph adjacency spectrum.

The research summary says a closed-form tensor spectrum is known in Cooper--Dutle and follow-up literature, but gives no exact formula or citation. The registry's open/OPEN-TRIAGE labels and that solved-literature report are retained separately.
\subsection{Short English statement}
Compute the Fano plane's homogeneous hypergraph adjacency-tensor spectrum, with the normalization and eigenvalue multiplicities fixed. The supplied summary reports a known closed form but provides no explicit formula or complete reference.
\subsection{Sources}
\begin{itemize}
\item Cooper, \emph{Combinatorial Problems I Like}, October 2020 snapshot, ``Graphs, Hypergraphs, Set Systems, Designs,'' source-order bullet 30, attributed to Clark, \url{https://people.math.sc.edu/cooper/combprob.html}; accessed 2026-07-29, HTML extraction.
\item Cooper--Dutle and follow-up hypergraph spectral literature, as the closed-form-spectrum attribution supplied by the research summary; no publication title, year, or formula is given.
\item Local registry \texttt{data/unsolvedmath/problems.json}, permanent UnsolvedMath ID $3100030$, source code catalogue entry, \emph{What is the (homogeneous adjacency) spectrum of the Fano plane}. Dataset attribution: Ulam AI and the UnsolvedMath Contributors.
\end{itemize}
\end{document}
